\chapter{Base de datos}

\section{Descripción general}
La base de datos relacional organiza, centraliza y relaciona de forma segura la información generada durante el ciclo de vida de una obra. Está diseñada para PostgreSQL y se define en el script \cod{db/esquema\_obras.sql}; los cambios sobre una base con datos se aplican con scripts incrementales (\cod{db/cu*.sql}) que pueden ejecutarse más de una vez y no borran información.

\subsection*{Decisiones de diseño}
\begin{itemize}
  \item \textbf{Identificadores UUID versión 4} como llave primaria en todas las tablas operativas. A diferencia de un contador, no revelan cuántos registros existen ni permiten adivinar identificadores ajenos (ataques IDOR). Los catálogos pequeños usan identificadores numéricos.
  \item \textbf{Integridad en la base de datos:} llaves foráneas, restricciones \cod{CHECK} y \cod{UNIQUE} y columnas calculadas garantizan la consistencia aunque falle la aplicación.
  \item \textbf{Disparadores} calculan los totales de las estimaciones y asignan el folio de la bitácora, de modo que no dependen de la aplicación.
  \item \textbf{Historial no editable:} las decisiones sobre reportes, estimaciones y notas se agregan como renglones nuevos; nunca se modifica una decisión anterior.
  \item \textbf{Esquema administrado por SQL:} los modelos de Django usan \cod{managed = False}, así que Django lee y escribe pero no altera la estructura.
  \item \textbf{Seguridad en Supabase:} la aplicación se conecta con un rol propio con permisos limitados y seguridad a nivel de renglón habilitada.
\end{itemize}

\section{Diagrama relacional}
\begin{sidewaysfigure}
\centering
\resizebox{\linewidth}{!}{%
\begin{tikzpicture}[
  tabla/.style={rectangle,rounded corners=2pt,draw=#1,fill=#1!8,line width=0.6pt,font=\ttfamily\scriptsize,minimum width=3.3cm,minimum height=0.62cm},
  fk/.style={-{Stealth[length=1.8mm]},draw=black!50,line width=0.45pt},
  grupo/.style={rectangle,rounded corners=6pt,draw=#1!40,dashed,line width=0.5pt,inner sep=7pt},
  etiq/.style={font=\sffamily\scriptsize\bfseries,text=#1}]
  % núcleo
  \node[tabla=guinda] (obras) at (-1.5,0) {obras};
  \node[tabla=guinda] (contratos) at (4,0) {contratos};
  \node[tabla=guinda] (anticipos) at (9.5,0) {anticipos};
  \node[tabla=guinda] (conceptos) at (4,-1.5) {catalogo\_conceptos};
  \node[tabla=guinda] (programa) at (-1.5,-1.5) {programa\_obra};
  % seguridad
  \node[tabla=azulinfo] (usuarios) at (-7,0) {usuarios};
  \node[tabla=azulinfo] (roles) at (-7,-3) {roles};
  \node[tabla=azulinfo] (upr) at (-7,-1.5) {usuarios\_proyectos\_roles};
  % operación
  \node[tabla=verdeok] (estim) at (9.5,-3.4) {estimaciones};
  \node[tabla=verdeok] (detalle) at (9.5,-5) {estimaciones\_detalles};
  \node[tabla=verdeok] (diarios) at (-1.5,-3.4) {reportes\_diarios};
  \node[tabla=verdeok] (rdn) at (1.6,-5) {reportes\_diarios\_notas};
  \node[tabla=verdeok] (mensuales) at (-7,-5) {reportes\_mensuales};
  \node[tabla=verdeok] (notas) at (4,-3.4) {notas\_bitacora};
  % auditoría
  \node[tabla=ambar] (vest) at (15,-3.4) {validaciones\_estimaciones};
  \node[tabla=ambar] (vrd) at (-1.6,-6.6) {validaciones\_reportes\_diarios};
  \node[tabla=ambar] (vrm) at (-7,-6.6) {validaciones\_reportes\_mensuales};
  \node[tabla=ambar] (firmas) at (7.6,-6.6) {firmas\_notas\_bitacora};
  \node[tabla=ambar] (docs) at (15,-5) {documentos};
  \node[tabla=ambar] (evid) at (3.4,-6.6) {evidencias};
  \node[tabla=ambar] (est) at (-2.5,-8.2) {catalogo\_estatus\_validacion};
  \node[tabla=ambar] (tipos) at (12.2,-6.6) {catalogo\_tipos\_nota};
  % relaciones (flecha hacia la tabla referenciada)
  \draw[fk] (contratos) -- (obras);
  \draw[fk] (anticipos) -- (contratos);
  \draw[fk] (conceptos) -- (obras);
  \draw[fk] (programa) -- (obras);
  \draw[fk] (upr) -- (usuarios);
  \draw[fk] (upr) -- (roles);
  \draw[fk] (upr.east) -- (obras.south west);
  \draw[fk] (estim) -- (contratos.south east);
  \draw[fk] (detalle) -- (estim);
  \draw[fk] (detalle.north west) -- (conceptos.south east);
  \draw[fk] (detalle.west) -- (notas.south east);
  \draw[fk,dashed] (diarios) -- (programa);
  \draw[fk] (rdn) -- (diarios);
  \draw[fk] (rdn) -- (notas);
  \draw[fk,dashed] (mensuales) -- (diarios.south west);
  \draw[fk] (notas.north west) -- (obras.south east);
  \draw[fk] (vest) -- (estim);
  \draw[fk] (vrd) -- (diarios);
  \draw[fk] (vrm) -- (mensuales);
  \draw[fk] (firmas) -- (notas);
  \draw[fk] (docs) -- (estim);
  \draw[fk] (docs.north) to[bend right=12] (obras.north east);
  \draw[fk] (evid) -- (diarios);
  \draw[fk] (evid) -- (mensuales);
  \draw[fk] (evid) -- (notas);
  \draw[fk] (notas) -- (tipos);
  % grupos
  \begin{scope}[on background layer]
    \node[grupo=guinda,fit=(obras)(contratos)(anticipos)(conceptos)(programa)] (gn) {};
    \node[grupo=azulinfo,fit=(usuarios)(upr)(roles)] (gs) {};
  \end{scope}
  \node[etiq=guinda,anchor=south west] at (gn.north west) {Núcleo};
  \node[etiq=azulinfo,anchor=south west] at (gs.north west) {Seguridad};
  \node[etiq=verdeok,anchor=west] at (13.2,-7.6) {\textcolor{verdeok}{$\blacksquare$} Operación};
  \node[etiq=ambar,anchor=west] at (13.2,-8.1) {\textcolor{ambar}{$\blacksquare$} Auditoría y archivos};
\end{tikzpicture}}
\caption[Diagrama relacional]{Diagrama relacional. Las flechas apuntan a la tabla referenciada. Por legibilidad se omiten las referencias a \cod{usuarios} (autor, validador, firmante) y a \cod{catalogo\_estatus\_validacion}, presentes en casi todas las tablas de operación y auditoría. La relación de \cod{reportes\_diarios} con \cod{programa\_obra} y de \cod{reportes\_mensuales} con \cod{reportes\_diarios} es lógica (por obra y fecha), no una llave foránea.}\label{fig:relacional}
\end{sidewaysfigure}

\section{Disparadores}
\begin{xltabular}{\linewidth}{@{}>{\ttfamily\small}p{5.4cm}>{\small}X@{}}
\caption{Disparadores de la base de datos}\label{tab:disparadores}\\
\toprule
\normalfont\sffamily\bfseries Disparador & \sffamily\bfseries Función\\\midrule\endfirsthead
\bottomrule\endfoot
trg\_detalle\_preparar & Antes de guardar un renglón de estimación: verifica que el concepto sea de la obra del contrato, toma el precio unitario del catálogo si no se indica, calcula el acumulado anterior y marca si es excedente.\\
trg\_detalle\_subtotal & Después de cambiar los renglones: recalcula el subtotal de la estimación.\\
trg\_estimacion\_totales & Antes de guardar la estimación: calcula retención por derechos, IVA y monto neto.\\
trg\_nota\_folio & Antes de insertar una nota: asigna el folio consecutivo por obra bloqueando la fila de la obra, sin choques entre usuarios.\\
\end{xltabular}

\chapter{Diccionario de datos}
El diccionario de datos se genera automáticamente desde \cod{db/esquema\_obras.sql} con el script \cod{generar\_diccionario.py}, por lo que coincide con la base de datos implementada. \emph{PK}: llave primaria; \emph{FK}: llave foránea; \emph{calculado}: columna generada por la base de datos.

% Generado por generar_diccionario.py a partir de db/esquema_obras.sql. No editar a mano.

\section{Módulo núcleo}
Alcance físico y financiero de la obra: la obra, su contrato y anticipos, el catálogo de conceptos y el programa de obra.

\tablabd{obras}{Obra pública: datos generales y fechas oficiales.}
\begin{xltabular}{\linewidth}{@{}>{\ttfamily\footnotesize}p{4.9cm}>{\ttfamily\footnotesize}p{2.4cm}>{\raggedright\arraybackslash\footnotesize}p{3.3cm}>{\footnotesize}X@{}}
\toprule\normalfont\sffamily\bfseries\footnotesize Campo & \normalfont\sffamily\bfseries\footnotesize Tipo & \sffamily\bfseries Restricciones & \sffamily\bfseries Descripción\\\midrule\endfirsthead
\toprule\normalfont\sffamily\bfseries\footnotesize Campo & \normalfont\sffamily\bfseries\footnotesize Tipo & \sffamily\bfseries Restricciones & \sffamily\bfseries Descripción\\\midrule\endhead
\bottomrule\endfoot
id & UUID & PK & Identificador único (UUID versión 4).\\
nombre & VARCHAR(255) & obligatorio & Nombre de la obra.\\
ubicacion & TEXT &  & Ubicación de la obra.\\
fecha\_inicio\_oficial & DATE &  & Inicio oficial; define el estado de la obra y el arranque del programa.\\
fecha\_fin\_oficial & DATE &  & Terminación oficial.\\
ultimo\_folio\_bitacora & INTEGER & obligatorio, por omisión 0 & Contador de folios de la bitácora (lo usa un disparador).\\
fecha\_registro & TIMESTAMPTZ & obligatorio, por omisión now() & Fecha y hora del registro.\\
\end{xltabular}

\tablabd{contratos}{Contrato de la obra (uno por obra) con montos, porcentajes y garantías.}
\begin{xltabular}{\linewidth}{@{}>{\ttfamily\footnotesize}p{4.9cm}>{\ttfamily\footnotesize}p{2.4cm}>{\raggedright\arraybackslash\footnotesize}p{3.3cm}>{\footnotesize}X@{}}
\toprule\normalfont\sffamily\bfseries\footnotesize Campo & \normalfont\sffamily\bfseries\footnotesize Tipo & \sffamily\bfseries Restricciones & \sffamily\bfseries Descripción\\\midrule\endfirsthead
\toprule\normalfont\sffamily\bfseries\footnotesize Campo & \normalfont\sffamily\bfseries\footnotesize Tipo & \sffamily\bfseries Restricciones & \sffamily\bfseries Descripción\\\midrule\endhead
\bottomrule\endfoot
id & UUID & PK & Identificador único (UUID versión 4).\\
obra\_id & UUID & FK → obras, obligatorio, único & Referencia a obras.\\
numero\_contrato & VARCHAR(100) & obligatorio, único & Número del contrato.\\
contratista & VARCHAR(255) & obligatorio & Empresa contratista.\\
fecha\_firma & DATE &  & Fecha de firma del contrato.\\
fecha\_inicio & DATE &  & Inicio de los trabajos según el contrato.\\
fecha\_fin & DATE &  & Terminación de los trabajos según el contrato.\\
monto\_total\_contratado & NUMERIC(15,2) & obligatorio & Monto contratado sin IVA.\\
porcentaje\_iva & NUMERIC(5,2) & obligatorio, por omisión 16 & Porcentaje de IVA.\\
porcentaje\_retencion & NUMERIC(5,2) & obligatorio, por omisión 0 & Retención por derechos (por ejemplo 5 al millar).\\
porcentaje\_anticipo & NUMERIC(5,2) & obligatorio, por omisión 0 & Porcentaje de anticipo (máximo 30 \%).\\
porcentaje\_retencion\_atraso & NUMERIC(5,2) & obligatorio, por omisión 5 & Retención económica por atraso (LOPSRM art. 46 Bis).\\
monto\_garantia\_cumplimiento & NUMERIC(15,2) &  & Garantía de cumplimiento; tope de penas convencionales.\\
numero\_fianza & VARCHAR(100) &  & Número de la fianza.\\
afianzadora & VARCHAR(255) &  & Institución afianzadora.\\
monto\_anticipo & NUMERIC(15,2) & calculado & Monto del anticipo (calculado).\\
fecha\_registro & TIMESTAMPTZ & obligatorio, por omisión now() & Fecha y hora del registro.\\
\end{xltabular}

\tablabd{anticipos}{Ministraciones del anticipo del contrato.}
\begin{xltabular}{\linewidth}{@{}>{\ttfamily\footnotesize}p{4.9cm}>{\ttfamily\footnotesize}p{2.4cm}>{\raggedright\arraybackslash\footnotesize}p{3.3cm}>{\footnotesize}X@{}}
\toprule\normalfont\sffamily\bfseries\footnotesize Campo & \normalfont\sffamily\bfseries\footnotesize Tipo & \sffamily\bfseries Restricciones & \sffamily\bfseries Descripción\\\midrule\endfirsthead
\toprule\normalfont\sffamily\bfseries\footnotesize Campo & \normalfont\sffamily\bfseries\footnotesize Tipo & \sffamily\bfseries Restricciones & \sffamily\bfseries Descripción\\\midrule\endhead
\bottomrule\endfoot
id & UUID & PK & Identificador único (UUID versión 4).\\
contrato\_id & UUID & FK → contratos, obligatorio & Referencia a contratos.\\
tipo & VARCHAR(30) & obligatorio, por omisión 'Inicio de obra', valores: Inicio de obra, Materiales, Otro & Inicio de obra, Materiales u Otro.\\
monto & NUMERIC(15,2) & obligatorio & Monto de la ministración.\\
fecha\_programada & DATE &  & Fecha programada de entrega.\\
fecha\_pago & DATE &  & Fecha en que se pagó.\\
estatus & VARCHAR(20) & obligatorio, por omisión 'Programado', valores: Programado, Pagado, Cancelado & Programado, Pagado o Cancelado.\\
observaciones & TEXT &  & Observaciones en texto libre.\\
fecha\_registro & TIMESTAMPTZ & obligatorio, por omisión now() & Fecha y hora del registro.\\
\end{xltabular}

\tablabd{catalogo\_conceptos}{Conceptos de trabajo de la obra con su volumen y precio unitario.}
\begin{xltabular}{\linewidth}{@{}>{\ttfamily\footnotesize}p{4.9cm}>{\ttfamily\footnotesize}p{2.4cm}>{\raggedright\arraybackslash\footnotesize}p{3.3cm}>{\footnotesize}X@{}}
\toprule\normalfont\sffamily\bfseries\footnotesize Campo & \normalfont\sffamily\bfseries\footnotesize Tipo & \sffamily\bfseries Restricciones & \sffamily\bfseries Descripción\\\midrule\endfirsthead
\toprule\normalfont\sffamily\bfseries\footnotesize Campo & \normalfont\sffamily\bfseries\footnotesize Tipo & \sffamily\bfseries Restricciones & \sffamily\bfseries Descripción\\\midrule\endhead
\bottomrule\endfoot
id & UUID & PK & Identificador único (UUID versión 4).\\
obra\_id & UUID & FK → obras, obligatorio & Referencia a obras.\\
capitulo & VARCHAR(100) &  & Capítulo o partida.\\
clave & VARCHAR(50) & obligatorio & Clave del concepto, única por obra.\\
descripcion & TEXT & obligatorio & Descripción en texto libre.\\
unidad\_medida & VARCHAR(20) & obligatorio & Unidad de medida.\\
volumen\_total\_contratado & NUMERIC(14,4) & obligatorio & Cantidad contratada.\\
precio\_unitario & NUMERIC(15,2) & obligatorio & Precio unitario.\\
importe\_contratado & NUMERIC(18,2) & calculado & Cantidad por precio unitario (calculado).\\
\end{xltabular}

\tablabd{programa\_obra}{Avance físico programado acumulado en fechas clave; entre ellas se interpola.}
\begin{xltabular}{\linewidth}{@{}>{\ttfamily\footnotesize}p{4.9cm}>{\ttfamily\footnotesize}p{2.4cm}>{\raggedright\arraybackslash\footnotesize}p{3.3cm}>{\footnotesize}X@{}}
\toprule\normalfont\sffamily\bfseries\footnotesize Campo & \normalfont\sffamily\bfseries\footnotesize Tipo & \sffamily\bfseries Restricciones & \sffamily\bfseries Descripción\\\midrule\endfirsthead
\toprule\normalfont\sffamily\bfseries\footnotesize Campo & \normalfont\sffamily\bfseries\footnotesize Tipo & \sffamily\bfseries Restricciones & \sffamily\bfseries Descripción\\\midrule\endhead
\bottomrule\endfoot
id & UUID & PK & Identificador único (UUID versión 4).\\
obra\_id & UUID & FK → obras, obligatorio & Referencia a obras.\\
fecha & DATE & obligatorio & Fecha clave del programa.\\
avance\_programado\_porcentaje & NUMERIC(5,2) & obligatorio & Avance físico programado acumulado a esa fecha.\\
fecha\_registro & TIMESTAMPTZ & obligatorio, por omisión now() & Fecha y hora del registro.\\
\end{xltabular}

\section{Módulo de seguridad}
Cuentas, roles y la asignación de roles por obra, que limita los permisos de cada persona a las obras en que participa.

\tablabd{usuarios}{Cuentas de usuario.}
\begin{xltabular}{\linewidth}{@{}>{\ttfamily\footnotesize}p{4.9cm}>{\ttfamily\footnotesize}p{2.4cm}>{\raggedright\arraybackslash\footnotesize}p{3.3cm}>{\footnotesize}X@{}}
\toprule\normalfont\sffamily\bfseries\footnotesize Campo & \normalfont\sffamily\bfseries\footnotesize Tipo & \sffamily\bfseries Restricciones & \sffamily\bfseries Descripción\\\midrule\endfirsthead
\toprule\normalfont\sffamily\bfseries\footnotesize Campo & \normalfont\sffamily\bfseries\footnotesize Tipo & \sffamily\bfseries Restricciones & \sffamily\bfseries Descripción\\\midrule\endhead
\bottomrule\endfoot
id & UUID & PK & Identificador único (UUID versión 4).\\
nombre & VARCHAR(100) & obligatorio & Nombre de la persona.\\
apellidos & VARCHAR(100) & obligatorio & Apellidos de la persona.\\
correo & CITEXT & obligatorio, único & Correo; sirve para iniciar sesión y recibir avisos.\\
password\_hash & VARCHAR(255) & obligatorio & Contraseña cifrada con función hash; nunca en texto claro.\\
activo & BOOLEAN & obligatorio, por omisión true & Una cuenta inactiva no puede iniciar sesión.\\
fecha\_creacion & TIMESTAMPTZ & obligatorio, por omisión now() & Fecha y hora de creación de la cuenta.\\
\end{xltabular}

\tablabd{roles}{Catálogo de roles: Superintendente Contratista, Supervisor, Residente, Directivo y Administrador.}
\begin{xltabular}{\linewidth}{@{}>{\ttfamily\footnotesize}p{4.9cm}>{\ttfamily\footnotesize}p{2.4cm}>{\raggedright\arraybackslash\footnotesize}p{3.3cm}>{\footnotesize}X@{}}
\toprule\normalfont\sffamily\bfseries\footnotesize Campo & \normalfont\sffamily\bfseries\footnotesize Tipo & \sffamily\bfseries Restricciones & \sffamily\bfseries Descripción\\\midrule\endfirsthead
\toprule\normalfont\sffamily\bfseries\footnotesize Campo & \normalfont\sffamily\bfseries\footnotesize Tipo & \sffamily\bfseries Restricciones & \sffamily\bfseries Descripción\\\midrule\endhead
\bottomrule\endfoot
id & UUID & PK & Identificador único (UUID versión 4).\\
nombre & VARCHAR(50) & obligatorio, único & Nombre del rol.\\
\end{xltabular}

\tablabd{usuarios\_proyectos\_roles}{Rol de cada usuario en cada obra (una persona puede tener varios).}
\begin{xltabular}{\linewidth}{@{}>{\ttfamily\footnotesize}p{4.9cm}>{\ttfamily\footnotesize}p{2.4cm}>{\raggedright\arraybackslash\footnotesize}p{3.3cm}>{\footnotesize}X@{}}
\toprule\normalfont\sffamily\bfseries\footnotesize Campo & \normalfont\sffamily\bfseries\footnotesize Tipo & \sffamily\bfseries Restricciones & \sffamily\bfseries Descripción\\\midrule\endfirsthead
\toprule\normalfont\sffamily\bfseries\footnotesize Campo & \normalfont\sffamily\bfseries\footnotesize Tipo & \sffamily\bfseries Restricciones & \sffamily\bfseries Descripción\\\midrule\endhead
\bottomrule\endfoot
usuario\_id & UUID & PK, FK → usuarios, obligatorio & Referencia a usuarios.\\
obra\_id & UUID & PK, FK → obras, obligatorio & Referencia a obras.\\
rol\_id & UUID & PK, FK → roles, obligatorio & Referencia a roles.\\
\end{xltabular}

\section{Módulo de operación}
Información que se genera en la obra: estimaciones, reportes diarios y mensuales y bitácora.

\tablabd{estimaciones}{Estimación de trabajos ejecutados en un periodo. Los totales los calculan disparadores.}
\begin{xltabular}{\linewidth}{@{}>{\ttfamily\footnotesize}p{4.9cm}>{\ttfamily\footnotesize}p{2.4cm}>{\raggedright\arraybackslash\footnotesize}p{3.3cm}>{\footnotesize}X@{}}
\toprule\normalfont\sffamily\bfseries\footnotesize Campo & \normalfont\sffamily\bfseries\footnotesize Tipo & \sffamily\bfseries Restricciones & \sffamily\bfseries Descripción\\\midrule\endfirsthead
\toprule\normalfont\sffamily\bfseries\footnotesize Campo & \normalfont\sffamily\bfseries\footnotesize Tipo & \sffamily\bfseries Restricciones & \sffamily\bfseries Descripción\\\midrule\endhead
\bottomrule\endfoot
id & UUID & PK & Identificador único (UUID versión 4).\\
contrato\_id & UUID & FK → contratos, obligatorio & Referencia a contratos.\\
superintendente\_id & UUID & FK → usuarios, obligatorio & Superintendente que la registró.\\
numero\_estimacion & INTEGER & obligatorio & Número consecutivo por contrato.\\
folio & VARCHAR(50) & obligatorio & Folio EST-001, EST-002… único por contrato.\\
tipo & VARCHAR(20) & obligatorio, por omisión 'Ordinaria', valores: Ordinaria, Extraordinaria, Finiquito & Tipo de estimación (Reglamento art. 130).\\
periodo\_inicio & DATE & obligatorio & Inicio del periodo estimado.\\
fecha\_corte & DATE & obligatorio & Fecha de corte del periodo (máximo un mes).\\
fecha\_presentacion & DATE &  & Fecha en que se presentó a revisión.\\
subtotal & NUMERIC(15,2) & obligatorio, por omisión 0 & Importe de los trabajos (suma de renglones).\\
amortizacion\_anticipo & NUMERIC(15,2) & obligatorio, por omisión 0 & Importe por porcentaje de anticipo (Reglamento art. 143).\\
retencion & NUMERIC(15,2) & obligatorio, por omisión 0 & Retención por derechos (calculada).\\
iva & NUMERIC(15,2) & obligatorio, por omisión 0 & IVA sobre el importe menos la amortización (calculado).\\
monto\_neto & NUMERIC(15,2) & obligatorio, por omisión 0 & Neto a pagar (calculado).\\
retencion\_atraso & NUMERIC(15,2) & obligatorio, por omisión 0 & Retención económica por atraso al corte.\\
devolucion\_retencion & NUMERIC(15,2) & obligatorio, por omisión 0 & Devolución de retenciones anteriores al regularizarse.\\
dias\_atraso & INTEGER &  & Días de atraso al corte respecto del programa.\\
atraso\_no\_imputable & BOOLEAN & obligatorio, por omisión false & El Superintendente justifica que el atraso no le es imputable.\\
justificacion\_atraso & TEXT &  & Justificación del atraso no imputable.\\
estatus\_id & SMALLINT & FK → catalogo\_estatus\_validacion, obligatorio, por omisión 1 & Estatus vigente (catálogo de estatus).\\
observaciones & TEXT &  & Observaciones en texto libre.\\
fecha\_envio & TIMESTAMPTZ &  & Fecha y hora en que se envió a revisión.\\
fecha\_modificacion & TIMESTAMPTZ &  & Fecha y hora de la última corrección.\\
fecha\_registro & TIMESTAMPTZ & obligatorio, por omisión now() & Fecha y hora del registro.\\
\end{xltabular}

\tablabd{estimaciones\_detalles}{Renglones de la estimación: cantidad ejecutada por concepto.}
\begin{xltabular}{\linewidth}{@{}>{\ttfamily\footnotesize}p{4.9cm}>{\ttfamily\footnotesize}p{2.4cm}>{\raggedright\arraybackslash\footnotesize}p{3.3cm}>{\footnotesize}X@{}}
\toprule\normalfont\sffamily\bfseries\footnotesize Campo & \normalfont\sffamily\bfseries\footnotesize Tipo & \sffamily\bfseries Restricciones & \sffamily\bfseries Descripción\\\midrule\endfirsthead
\toprule\normalfont\sffamily\bfseries\footnotesize Campo & \normalfont\sffamily\bfseries\footnotesize Tipo & \sffamily\bfseries Restricciones & \sffamily\bfseries Descripción\\\midrule\endhead
\bottomrule\endfoot
id & UUID & PK & Identificador único (UUID versión 4).\\
estimacion\_id & UUID & FK → estimaciones, obligatorio & Referencia a estimaciones.\\
concepto\_id & UUID & FK → catalogo\_conceptos, obligatorio & Referencia a catalogo\_conceptos.\\
cantidad\_ejecutada & NUMERIC(14,4) & obligatorio & Cantidad ejecutada en el periodo.\\
cantidad\_acumulada\_anterior & NUMERIC(14,4) & obligatorio, por omisión 0 & Cantidad estimada en estimaciones anteriores (disparador).\\
cantidad\_acumulada & NUMERIC(14,4) & calculado & Acumulado anterior más el periodo (calculado).\\
precio\_unitario & NUMERIC(15,2) & obligatorio & Precio unitario; si va vacío se toma el del catálogo (disparador).\\
importe & NUMERIC(15,2) & calculado & Cantidad por precio unitario (calculado).\\
es\_excedente & BOOLEAN & obligatorio, por omisión false & Verdadero si el acumulado rebasa lo contratado (disparador).\\
numero\_generador & VARCHAR(50) &  & Número del generador que respalda la cantidad.\\
observaciones & TEXT &  & Observaciones en texto libre.\\
nota\_bitacora\_id & UUID & FK → notas\_bitacora & Nota de bitácora firmada que autorizó el excedente (Reglamento art. 118).\\
\end{xltabular}

\tablabd{reportes\_diarios}{Reporte diario del Supervisor con el avance físico real.}
\begin{xltabular}{\linewidth}{@{}>{\ttfamily\footnotesize}p{4.9cm}>{\ttfamily\footnotesize}p{2.4cm}>{\raggedright\arraybackslash\footnotesize}p{3.3cm}>{\footnotesize}X@{}}
\toprule\normalfont\sffamily\bfseries\footnotesize Campo & \normalfont\sffamily\bfseries\footnotesize Tipo & \sffamily\bfseries Restricciones & \sffamily\bfseries Descripción\\\midrule\endfirsthead
\toprule\normalfont\sffamily\bfseries\footnotesize Campo & \normalfont\sffamily\bfseries\footnotesize Tipo & \sffamily\bfseries Restricciones & \sffamily\bfseries Descripción\\\midrule\endhead
\bottomrule\endfoot
id & UUID & PK & Identificador único (UUID versión 4).\\
obra\_id & UUID & FK → obras, obligatorio & Referencia a obras.\\
supervisor\_id & UUID & FK → usuarios, obligatorio & Supervisor que elaboró el reporte.\\
fecha\_reporte & DATE & obligatorio & Fecha del reporte (uno por obra y fecha).\\
jornada\_inicio & TIME &  & Hora de inicio de la jornada.\\
jornada\_fin & TIME &  & Hora de fin de la jornada.\\
clima & VARCHAR(20) & valores: Soleado, Nublado, Lluvia, Tormenta, Viento fuerte, Otro & Condición del clima.\\
temperatura\_c & NUMERIC(4,1) &  & Temperatura en grados Celsius.\\
actividades\_realizadas & TEXT & obligatorio & Actividades realizadas en el día.\\
incidencias & TEXT &  & Incidencias del día.\\
incidentes\_seguridad & TEXT &  & Incidentes de seguridad del día.\\
avance\_fisico\_real\_porcentaje & NUMERIC(5,2) & obligatorio & Avance físico real acumulado (0 a 100).\\
observaciones & TEXT &  & Observaciones en texto libre.\\
atraso\_causa & TEXT &  & Causa del atraso crítico (1 día o más).\\
atraso\_solicitud & VARCHAR(20) &  & Solicitud ante el atraso crítico.\\
estatus\_id & SMALLINT & FK → catalogo\_estatus\_validacion, obligatorio, por omisión 1 & Estatus vigente (catálogo de estatus).\\
fecha\_registro & TIMESTAMPTZ & obligatorio, por omisión now() & Fecha y hora del registro.\\
\end{xltabular}

\tablabd{reportes\_diarios\_notas}{Notas de bitácora vinculadas a un reporte diario.}
\begin{xltabular}{\linewidth}{@{}>{\ttfamily\footnotesize}p{4.9cm}>{\ttfamily\footnotesize}p{2.4cm}>{\raggedright\arraybackslash\footnotesize}p{3.3cm}>{\footnotesize}X@{}}
\toprule\normalfont\sffamily\bfseries\footnotesize Campo & \normalfont\sffamily\bfseries\footnotesize Tipo & \sffamily\bfseries Restricciones & \sffamily\bfseries Descripción\\\midrule\endfirsthead
\toprule\normalfont\sffamily\bfseries\footnotesize Campo & \normalfont\sffamily\bfseries\footnotesize Tipo & \sffamily\bfseries Restricciones & \sffamily\bfseries Descripción\\\midrule\endhead
\bottomrule\endfoot
id & UUID & PK & Identificador único (UUID versión 4).\\
reporte\_diario\_id & UUID & FK → reportes\_diarios, obligatorio & Referencia a reportes\_diarios.\\
nota\_bitacora\_id & UUID & FK → notas\_bitacora, obligatorio & Referencia a notas\_bitacora.\\
\end{xltabular}

\tablabd{reportes\_mensuales}{Reporte mensual que consolida los reportes diarios del mes.}
\begin{xltabular}{\linewidth}{@{}>{\ttfamily\footnotesize}p{4.9cm}>{\ttfamily\footnotesize}p{2.4cm}>{\raggedright\arraybackslash\footnotesize}p{3.3cm}>{\footnotesize}X@{}}
\toprule\normalfont\sffamily\bfseries\footnotesize Campo & \normalfont\sffamily\bfseries\footnotesize Tipo & \sffamily\bfseries Restricciones & \sffamily\bfseries Descripción\\\midrule\endfirsthead
\toprule\normalfont\sffamily\bfseries\footnotesize Campo & \normalfont\sffamily\bfseries\footnotesize Tipo & \sffamily\bfseries Restricciones & \sffamily\bfseries Descripción\\\midrule\endhead
\bottomrule\endfoot
id & UUID & PK & Identificador único (UUID versión 4).\\
obra\_id & UUID & FK → obras, obligatorio & Referencia a obras.\\
supervisor\_id & UUID & FK → usuarios, obligatorio & Supervisor que elaboró el reporte.\\
anio & SMALLINT & obligatorio & Año del periodo.\\
mes & SMALLINT & obligatorio & Mes del periodo (1 a 12).\\
fecha\_corte & DATE & obligatorio & Fecha de corte, dentro del mes reportado.\\
resumen & TEXT &  & Resumen del mes redactado por el Supervisor.\\
observaciones & TEXT &  & Observaciones en texto libre.\\
estatus\_id & SMALLINT & FK → catalogo\_estatus\_validacion, obligatorio, por omisión 1 & Estatus vigente (catálogo de estatus).\\
consolidado & JSONB &  & Cifras del mes congeladas al aprobarse (JSON).\\
fecha\_registro & TIMESTAMPTZ & obligatorio, por omisión now() & Fecha y hora del registro.\\
fecha\_modificacion & TIMESTAMPTZ &  & Fecha y hora de la última corrección.\\
\end{xltabular}

\tablabd{notas\_bitacora}{Notas de la bitácora de obra con folio consecutivo por obra.}
\begin{xltabular}{\linewidth}{@{}>{\ttfamily\footnotesize}p{4.9cm}>{\ttfamily\footnotesize}p{2.4cm}>{\raggedright\arraybackslash\footnotesize}p{3.3cm}>{\footnotesize}X@{}}
\toprule\normalfont\sffamily\bfseries\footnotesize Campo & \normalfont\sffamily\bfseries\footnotesize Tipo & \sffamily\bfseries Restricciones & \sffamily\bfseries Descripción\\\midrule\endfirsthead
\toprule\normalfont\sffamily\bfseries\footnotesize Campo & \normalfont\sffamily\bfseries\footnotesize Tipo & \sffamily\bfseries Restricciones & \sffamily\bfseries Descripción\\\midrule\endhead
\bottomrule\endfoot
id & UUID & PK & Identificador único (UUID versión 4).\\
obra\_id & UUID & FK → obras, obligatorio & Referencia a obras.\\
folio & INTEGER & obligatorio, por omisión 0 & Folio consecutivo por obra; nunca se reutiliza.\\
tipo\_id & SMALLINT & FK → catalogo\_tipos\_nota, obligatorio, por omisión 7 & Tipo de nota (catálogo).\\
usuario\_creador\_id & UUID & FK → usuarios, obligatorio & Autor de la nota.\\
fecha\_nota & TIMESTAMPTZ & obligatorio, por omisión now() & Fecha y hora de la nota.\\
fecha\_modificacion & TIMESTAMPTZ &  & Fecha y hora de la última corrección.\\
descripcion & TEXT & obligatorio & Descripción en texto libre.\\
nota\_relacionada\_id & UUID & FK → notas\_bitacora & Nota a la que responde.\\
estado\_firma & VARCHAR(20) & obligatorio, por omisión 'Pendiente', valores: Pendiente, Firmada, Rechazada & Pendiente, Firmada o Rechazada.\\
\end{xltabular}

\section{Módulo de auditoría y archivos}
Catálogos de estatus, historial no editable de validaciones y firmas, y repositorio de documentos y evidencias.

\tablabd{catalogo\_estatus\_validacion}{Estatus: Pendiente, Aprobado, Rechazado, En revisión, Borrador y Con observaciones.}
\begin{xltabular}{\linewidth}{@{}>{\ttfamily\footnotesize}p{4.9cm}>{\ttfamily\footnotesize}p{2.4cm}>{\raggedright\arraybackslash\footnotesize}p{3.3cm}>{\footnotesize}X@{}}
\toprule\normalfont\sffamily\bfseries\footnotesize Campo & \normalfont\sffamily\bfseries\footnotesize Tipo & \sffamily\bfseries Restricciones & \sffamily\bfseries Descripción\\\midrule\endfirsthead
\toprule\normalfont\sffamily\bfseries\footnotesize Campo & \normalfont\sffamily\bfseries\footnotesize Tipo & \sffamily\bfseries Restricciones & \sffamily\bfseries Descripción\\\midrule\endhead
\bottomrule\endfoot
id & SMALLINT & PK, calculado, por omisión AS & Identificador del catálogo.\\
nombre & VARCHAR(50) & obligatorio, único & Nombre del estatus.\\
descripcion & TEXT &  & Descripción en texto libre.\\
\end{xltabular}

\tablabd{catalogo\_tipos\_nota}{Tipos de nota de bitácora: Apertura, Instrucción, Solicitud, Avance, Incidencia, Cierre y Otra.}
\begin{xltabular}{\linewidth}{@{}>{\ttfamily\footnotesize}p{4.9cm}>{\ttfamily\footnotesize}p{2.4cm}>{\raggedright\arraybackslash\footnotesize}p{3.3cm}>{\footnotesize}X@{}}
\toprule\normalfont\sffamily\bfseries\footnotesize Campo & \normalfont\sffamily\bfseries\footnotesize Tipo & \sffamily\bfseries Restricciones & \sffamily\bfseries Descripción\\\midrule\endfirsthead
\toprule\normalfont\sffamily\bfseries\footnotesize Campo & \normalfont\sffamily\bfseries\footnotesize Tipo & \sffamily\bfseries Restricciones & \sffamily\bfseries Descripción\\\midrule\endhead
\bottomrule\endfoot
id & SMALLINT & PK, calculado, por omisión AS & Identificador del catálogo.\\
nombre & VARCHAR(50) & obligatorio, único & Nombre del tipo de nota.\\
descripcion & TEXT &  & Descripción en texto libre.\\
\end{xltabular}

\tablabd{validaciones\_reportes\_diarios}{Historial de decisiones sobre los reportes diarios.}
\begin{xltabular}{\linewidth}{@{}>{\ttfamily\footnotesize}p{4.9cm}>{\ttfamily\footnotesize}p{2.4cm}>{\raggedright\arraybackslash\footnotesize}p{3.3cm}>{\footnotesize}X@{}}
\toprule\normalfont\sffamily\bfseries\footnotesize Campo & \normalfont\sffamily\bfseries\footnotesize Tipo & \sffamily\bfseries Restricciones & \sffamily\bfseries Descripción\\\midrule\endfirsthead
\toprule\normalfont\sffamily\bfseries\footnotesize Campo & \normalfont\sffamily\bfseries\footnotesize Tipo & \sffamily\bfseries Restricciones & \sffamily\bfseries Descripción\\\midrule\endhead
\bottomrule\endfoot
id & UUID & PK & Identificador único (UUID versión 4).\\
reporte\_diario\_id & UUID & FK → reportes\_diarios, obligatorio & Referencia a reportes\_diarios.\\
usuario\_validador\_id & UUID & FK → usuarios, obligatorio & Usuario que registró la decisión.\\
rol\_validador & VARCHAR(40) & valores: Residente, Superintendente Contratista & Rol con el que se decidió.\\
estatus\_id & SMALLINT & FK → catalogo\_estatus\_validacion, obligatorio & Estatus vigente (catálogo de estatus).\\
fecha\_validacion & TIMESTAMPTZ & obligatorio, por omisión now() & Fecha y hora de la decisión.\\
comentarios & TEXT &  & Comentarios o motivo de la decisión.\\
\end{xltabular}

\tablabd{validaciones\_reportes\_mensuales}{Historial de decisiones sobre los reportes mensuales.}
\begin{xltabular}{\linewidth}{@{}>{\ttfamily\footnotesize}p{4.9cm}>{\ttfamily\footnotesize}p{2.4cm}>{\raggedright\arraybackslash\footnotesize}p{3.3cm}>{\footnotesize}X@{}}
\toprule\normalfont\sffamily\bfseries\footnotesize Campo & \normalfont\sffamily\bfseries\footnotesize Tipo & \sffamily\bfseries Restricciones & \sffamily\bfseries Descripción\\\midrule\endfirsthead
\toprule\normalfont\sffamily\bfseries\footnotesize Campo & \normalfont\sffamily\bfseries\footnotesize Tipo & \sffamily\bfseries Restricciones & \sffamily\bfseries Descripción\\\midrule\endhead
\bottomrule\endfoot
id & UUID & PK & Identificador único (UUID versión 4).\\
reporte\_mensual\_id & UUID & FK → reportes\_mensuales, obligatorio & Referencia a reportes\_mensuales.\\
usuario\_validador\_id & UUID & FK → usuarios, obligatorio & Usuario que registró la decisión.\\
rol\_validador & VARCHAR(40) & obligatorio, valores: Residente, Superintendente Contratista & Rol con el que se decidió.\\
estatus\_id & SMALLINT & FK → catalogo\_estatus\_validacion, obligatorio & Estatus vigente (catálogo de estatus).\\
fecha\_validacion & TIMESTAMPTZ & obligatorio, por omisión now() & Fecha y hora de la decisión.\\
comentarios & TEXT &  & Comentarios o motivo de la decisión.\\
\end{xltabular}

\tablabd{validaciones\_estimaciones}{Historial de decisiones del Supervisor y del Residente sobre las estimaciones.}
\begin{xltabular}{\linewidth}{@{}>{\ttfamily\footnotesize}p{4.9cm}>{\ttfamily\footnotesize}p{2.4cm}>{\raggedright\arraybackslash\footnotesize}p{3.3cm}>{\footnotesize}X@{}}
\toprule\normalfont\sffamily\bfseries\footnotesize Campo & \normalfont\sffamily\bfseries\footnotesize Tipo & \sffamily\bfseries Restricciones & \sffamily\bfseries Descripción\\\midrule\endfirsthead
\toprule\normalfont\sffamily\bfseries\footnotesize Campo & \normalfont\sffamily\bfseries\footnotesize Tipo & \sffamily\bfseries Restricciones & \sffamily\bfseries Descripción\\\midrule\endhead
\bottomrule\endfoot
id & UUID & PK & Identificador único (UUID versión 4).\\
estimacion\_id & UUID & FK → estimaciones, obligatorio & Referencia a estimaciones.\\
usuario\_validador\_id & UUID & FK → usuarios, obligatorio & Usuario que registró la decisión.\\
estatus\_id & SMALLINT & FK → catalogo\_estatus\_validacion, obligatorio & Estatus vigente (catálogo de estatus).\\
fecha\_validacion & TIMESTAMPTZ & obligatorio, por omisión now() & Fecha y hora de la decisión.\\
comentarios & TEXT &  & Comentarios o motivo de la decisión.\\
rol\_validador & VARCHAR(40) & valores: Supervisor, Residente & Rol con el que se decidió.\\
\end{xltabular}

\tablabd{firmas\_notas\_bitacora}{Historial de firmas de cada rol sobre las notas de bitácora.}
\begin{xltabular}{\linewidth}{@{}>{\ttfamily\footnotesize}p{4.9cm}>{\ttfamily\footnotesize}p{2.4cm}>{\raggedright\arraybackslash\footnotesize}p{3.3cm}>{\footnotesize}X@{}}
\toprule\normalfont\sffamily\bfseries\footnotesize Campo & \normalfont\sffamily\bfseries\footnotesize Tipo & \sffamily\bfseries Restricciones & \sffamily\bfseries Descripción\\\midrule\endfirsthead
\toprule\normalfont\sffamily\bfseries\footnotesize Campo & \normalfont\sffamily\bfseries\footnotesize Tipo & \sffamily\bfseries Restricciones & \sffamily\bfseries Descripción\\\midrule\endhead
\bottomrule\endfoot
id & UUID & PK & Identificador único (UUID versión 4).\\
nota\_bitacora\_id & UUID & FK → notas\_bitacora, obligatorio & Referencia a notas\_bitacora.\\
usuario\_firmante\_id & UUID & FK → usuarios, obligatorio & Usuario que firmó.\\
rol\_firmante & VARCHAR(40) & obligatorio, valores: Residente, Supervisor, Superintendente Contratista & Rol con el que se firmó.\\
decision & VARCHAR(20) & obligatorio, valores: Firmada, Rechazada, Pendiente & Firmada, Rechazada o Pendiente (reinicio).\\
fecha\_firma & TIMESTAMPTZ & obligatorio, por omisión now() & Fecha y hora de la firma.\\
comentarios & TEXT &  & Comentarios o motivo de la decisión.\\
\end{xltabular}

\tablabd{documentos}{Repositorio de documentos de la obra y soporte de estimaciones.}
\begin{xltabular}{\linewidth}{@{}>{\ttfamily\footnotesize}p{4.9cm}>{\ttfamily\footnotesize}p{2.4cm}>{\raggedright\arraybackslash\footnotesize}p{3.3cm}>{\footnotesize}X@{}}
\toprule\normalfont\sffamily\bfseries\footnotesize Campo & \normalfont\sffamily\bfseries\footnotesize Tipo & \sffamily\bfseries Restricciones & \sffamily\bfseries Descripción\\\midrule\endfirsthead
\toprule\normalfont\sffamily\bfseries\footnotesize Campo & \normalfont\sffamily\bfseries\footnotesize Tipo & \sffamily\bfseries Restricciones & \sffamily\bfseries Descripción\\\midrule\endhead
\bottomrule\endfoot
id & UUID & PK & Identificador único (UUID versión 4).\\
obra\_id & UUID & FK → obras, obligatorio & Referencia a obras.\\
usuario\_subida\_id & UUID & FK → usuarios, obligatorio & Usuario que lo subió.\\
nombre\_archivo & VARCHAR(255) & obligatorio & Nombre del documento.\\
ruta\_almacenamiento & VARCHAR(500) & obligatorio & Ruta del archivo en Supabase Storage.\\
tipo\_archivo & VARCHAR(50) &  & Formato del archivo (PDF, DOCX, XLSX, imagen o DWG).\\
descripcion & TEXT &  & Descripción en texto libre.\\
fecha\_subida & TIMESTAMPTZ & obligatorio, por omisión now() & Fecha y hora de la carga.\\
categoria & VARCHAR(40) & obligatorio, por omisión 'Otro' & Tipo de documento.\\
estimacion\_id & UUID & FK → estimaciones & Estimación de la que es soporte (opcional).\\
\end{xltabular}

\tablabd{evidencias}{Fotografías y PDF de un reporte diario, un reporte mensual o una nota de bitácora.}
\begin{xltabular}{\linewidth}{@{}>{\ttfamily\footnotesize}p{4.9cm}>{\ttfamily\footnotesize}p{2.4cm}>{\raggedright\arraybackslash\footnotesize}p{3.3cm}>{\footnotesize}X@{}}
\toprule\normalfont\sffamily\bfseries\footnotesize Campo & \normalfont\sffamily\bfseries\footnotesize Tipo & \sffamily\bfseries Restricciones & \sffamily\bfseries Descripción\\\midrule\endfirsthead
\toprule\normalfont\sffamily\bfseries\footnotesize Campo & \normalfont\sffamily\bfseries\footnotesize Tipo & \sffamily\bfseries Restricciones & \sffamily\bfseries Descripción\\\midrule\endhead
\bottomrule\endfoot
id & UUID & PK & Identificador único (UUID versión 4).\\
ruta & VARCHAR(500) & obligatorio & Ruta del archivo en Supabase Storage.\\
tipo & VARCHAR(50) &  & Tipo de contenido (imagen o PDF).\\
descripcion & TEXT &  & Descripción en texto libre.\\
usuario\_subida\_id & UUID & FK → usuarios, obligatorio & Usuario que la subió.\\
fecha\_subida & TIMESTAMPTZ & obligatorio, por omisión now() & Fecha y hora de la carga.\\
reporte\_diario\_id & UUID & FK → reportes\_diarios & Referencia a reportes\_diarios.\\
reporte\_mensual\_id & UUID & FK → reportes\_mensuales & Referencia a reportes\_mensuales.\\
nota\_bitacora\_id & UUID & FK → notas\_bitacora & Referencia a notas\_bitacora.\\
\end{xltabular}

